\section{Conclusion}

\subsection{Selected symmetry and full symmetry}

The canonical-cover construction selects a quotient frame

\[
\mathcal P_0
\]

and exposes its stabilizer

\[
A_{\mathcal P_0}
\cong
S_5\times C_2.
\]

The unmarked graph has the larger symmetry

\[
\Aut(X)
\cong
S_5\times S_3.
\]

The enlargement is explained by the homogeneous-space realization. The
complementary \(S_3\) is the centralizer of the transitive \(S_5\)-action,
and its three involutions are the three block swaps in the quotient-frame
orbit.

Thus the visible \(C_2\) is not lost. It is recognized as one selected
transposition inside a larger intrinsic \(S_3\).

\subsection{Two complementary descriptions}

The graph admits two principal constructions:

\[
X=\KC(L(P))
\]

and

\[
X=
\text{the connected quartic self-paired orbital graph on }
S_5/V_{4,\mathrm{even}}.
\]

The first presentation emphasizes fibers, deck transformations, and
quotients. The second emphasizes orbitals, normalizers, centralizers, and
weak automorphisms of the Schurian configuration.

Their coincidence explains the automorphism group. Neither presentation
alone displays the complete structure as efficiently as the two together.

The invariant cohomology calculation gives a third perspective. The graph
is the all-one member

\[
\mathcal C_{11}
\]

of the four-element fixed cover square

\[
H^1(L(P),\Ftwo)^{S_5}\cong\Ftwo^2.
\]

\subsection{Open problems}

Three questions remain immediate.

\begin{enumerate}

\item
Classify all regular \(L(P)\)-quotient frames of \(X\), or prove that the
three-member orbit

\[
\frorbit
\]

is complete.

\item
Classify the four quartic self-paired double cosets of
\(V_{4,\mathrm{even}}\) in \(S_5\) conceptually, without exhaustive
enumeration.

\item
Derive the two-dimensional fixed cohomology space and its
triangle-pentagon coordinate basis representation-theoretically.

\end{enumerate}

The present paper establishes the orbit theorem rather than the first
classification, the certified orbital census rather than the second, and
the exact fixed-space computation rather than the third.
